\documentclass[preview]{standalone}

\usepackage{amsmath}
\usepackage{amssymb}
\usepackage{parskip}
\usepackage{fullpage}
\usepackage{hyperref}
\usepackage{stellar}
\usepackage{definitions}
\usepackage{bettelini}

\begin{document}

\id{taylor-polynomials}
\genpage

\section{Prerequisites}

\begin{snippetproposition}{power-rule-nth-derivative}{\(n\)-th derivative of \(x^n\)}
    The \(n\)-th derivative of \(x^n\), for \(n \in {\naturalnumbers}^*\),
    is given by
    \[
        \derivativeD^n x^n = n\factorial
    \]
\end{snippetproposition}

\begin{snippetproof}{power-rule-nth-derivative-proof}{power-rule-nth-derivative}{\(n\)-th derivative of \(x^n\)}
    \begin{align*}
        \frac{d^n}{dx^n}\left(x^n\right)
        &=\frac{d^{(n-1)}}{dx^{(n-1)}}\left(nx^{(n-1)}\right)\\
        &=\frac{d^{(n-2)}}{dx^{(n-2)}}\left(n(n-1)x^{(n-2)}\right)\\
        &=\cdots\\
        &=\prod_{i=1}^n i=n\factorial.
    \end{align*}
\end{snippetproof}

\begin{snippetlemma}{centered-polynomial-derivative}{Centered polynomial derivative}
    Let
    \[
        P(x) = \sum_{k=0}^n a_k{(x-x_0)}^k
    \]
    be a \polynomial of degree at most \(n\), written in powers of \(x-x_0\).
    Then, for every \(m \in \{0,\ldots,n\}\),
    \[
        P^{(m)}(x_0)=m\factorial a_m.
    \]
\end{snippetlemma}

\section{Introduction}

\begin{snippet}{taylor-series-introduction-derivation}
    We want to construct a power series centered around the point \(x=a\),
    so the variable will be \(x-a\).
    \begin{align*}
        \sum_{n=0}^{\infty}c_n{(x-a)}^n
    \end{align*}
    The goal is to find the coefficients \(c_n\).
    We first notice that \(f(a)=c_0\), which is the only coefficient that does not multiply a variable.
    If we take the derivative, the coefficient \(c_1\) will lose its variable
    \begin{align*}
        c_1+2c_2(x-a)+3c_3{(x-a)}^2+\cdots
    \end{align*}
    Now we have \(f'(a)=c_1\).
    \\
    We take the derivative of the polynomial again
    \begin{align*}
        2c_2+6c_3(x-a)+\cdots
    \end{align*}
    This time we have
    \begin{align*}
        f''(a)&=2c_2\\
        c_2&=\frac{f''(a)}{2}
    \end{align*}
    And again
    \begin{align*}
        f'''(a)&=6c_3\\
        c_3&=\frac{f'''(a)}{6}
    \end{align*}

    More generally, to extract the coefficient \(c_n\), we take the \(n\)-th derivative of the function. \\
    This brings us to
    \begin{align*}
        c_n=\frac{f^{(n)}(a)}{n\factorial}
    \end{align*}
    and finally to
    \begin{align*}
        \sum_{n=0}^{\infty}\frac{{(x-a)}^n f^{(n)}(a)}{n\factorial}
    \end{align*}
    This series need not converge, and even when it does, its sum need not coincide with the function.
\end{snippet}

\section{Taylor Polynomial}

\begin{snippetdefinition}{taylor-polynomial-definition}{Taylor polynomial}
    Let \(n\in\naturalnumbers\), and let \(f\colon I \fromto \realnumbers\) be
    \(n\) times differentiable at an interior point \(x_0\) of \(I\).
    The \emph{Taylor polynomial of order \(n\) centered at \(x_0\)}
    is defined as
    \[
        P_n(f, x_0, x) \triangleq \sum_{k=0}^n
        \frac{f^{(k)}(x_0) \cdot {(x-x_0)}^k}{k!}
    \]
\end{snippetdefinition}

\begin{snippetdefinition}{maclaurin-polynomial-definition}{Maclaurin polynomial}
    A \emph{Maclaurin polynomial} is a Taylor polynomial centered at \(x_0 = 0\).
\end{snippetdefinition}

\begin{snippetproposition}{taylor-polynomial-unique-property}{Characterization of the Taylor polynomial}
    Let \(n\in\naturalnumbers\), and let \(f\colon I\to\realnumbers\) be
    \(n\) times differentiable at an interior point \(x_0\) of \(I\).
    The Taylor polynomial of \(f\) of order \(n\), centered at \(x_0\), is the
    unique \polynomial \(P\) of degree at most \(n\) such that,
    for every \(k \in \{0,\ldots,n\}\),
    \[
        P^{(k)}(x_0)=f^{(k)}(x_0).
    \]
\end{snippetproposition}

\begin{plainhtml}
    The next theorem extends the local linear approximation of differentiable functions
    to approximation by higher-degree polynomials.
\end{plainhtml}

\begin{snippettheorem}{taylor-theorem}{Taylor's theorem}
    Let \(n\in\naturalnumbers^*\), and let \(f\colon I \to \realnumbers\) be
    \(n\) times differentiable at an interior point \(x_0\) of the interval \(I\).
    Let
    \[
        P_n(x) = \sum_{k=0}^n \frac{f^{(k)}(x_0){(x-x_0)}^k}{k!}
    \]
    be its Taylor polynomial of order \(n\), centered at \(x_0\).
    Then,
    \begin{enumerate}
        \item \emph{Taylor's formula with the Peano form of the remainder:}
        \[
            f(x)=P_n(x)+r_n(x),
            \qquad
            r_n(x)=\littleO\!\left({(x-x_0)}^n\right)
            \quad\text{as }x\to x_0.
        \]
        \item \emph{Taylor's formula with the Lagrange form of the remainder:}
        if \(f\) is \((n+1)\) times differentiable on \(I\), then, for every
        \(x\in I\setminus\{x_0\}\), there exists a point \(c_x\) strictly between
        \(x_0\) and \(x\) such that
        \[
            f(x)=P_n(x)+r_n(x),
            \qquad
            r_n(x)=\frac{f^{(n+1)}(c_x)}{(n+1)!}(x-x_0)^{n+1}.
        \]
        In general, \(c_x\) depends on both \(x\) and \(n\).
    \end{enumerate}
\end{snippettheorem}

\begin{snippetproof}{taylor-theorem-proof}{taylor-theorem}{Taylor's theorem}
    \begin{enumerate}
        \item We prove the Peano formula by induction on \(n\). For \(n=1\),
        differentiability at \(x_0\) gives
        \[
            f(x)=f(x_0)+f'(x_0)(x-x_0)+\littleO(x-x_0).
        \]
        Suppose the result holds at order \(n-1\). Applying the induction
        hypothesis to \(f'\) yields
        \[
            f'(x)-P_n'(x)=\littleO\!\left({(x-x_0)}^{n-1}\right).
        \]
        Set \(g(x)=f(x)-P_n(x)\). Since \(g(x_0)=0\), the mean value theorem
        gives, for every \(x\neq x_0\) sufficiently close to \(x_0\), a point
        \(\xi_x\) between \(x_0\) and \(x\) such that
        \[
            g(x)=g'(\xi_x)(x-x_0).
        \]
        As \(x\to x_0\), we also have \(\xi_x\to x_0\), and therefore
        \[
            g'(\xi_x)=\littleO\!\left({(\xi_x-x_0)}^{n-1}\right)
            =\littleO\!\left({(x-x_0)}^{n-1}\right).
        \]
        Hence \(g(x)=\littleO\!\left({(x-x_0)}^n\right)\), which proves the
        Peano form of the remainder.

        \item Fix \(x\in I\setminus\{x_0\}\), and write
        \[
            R=f(x)-P_n(x).
        \]
        Define, on the interval with endpoints \(x_0\) and \(x\),
        \[
            \Phi(t)=f(t)-P_n(t)
            -R\left(\frac{t-x_0}{x-x_0}\right)^{n+1}.
        \]
        We have \(\Phi(x_0)=\Phi(x)=0\), and the defining property of the
        Taylor polynomial gives
        \[
            \Phi^{(k)}(x_0)=0
            \qquad\text{for every }k\in\{1,\ldots,n\}.
        \]
        Repeated application of Rolle's theorem therefore yields a point \(c_x\)
        strictly between \(x_0\) and \(x\) such that
        \(\Phi^{(n+1)}(c_x)=0\). Since \(P_n^{(n+1)}=0\),
        \[
            0=f^{(n+1)}(c_x)-R\frac{(n+1)!}{(x-x_0)^{n+1}}.
        \]
        Solving for \(R\) gives
        \[
            R=\frac{f^{(n+1)}(c_x)}{(n+1)!}(x-x_0)^{n+1},
        \]
        which is the Lagrange form of the remainder.
    \end{enumerate}
\end{snippetproof}

\begin{snippettheorem}{taylor-uniqueness-theorem}{Uniqueness of the Taylor polynomial}
    Let \(n\in\naturalnumbers^*\), and let \(f\colon I \to \realnumbers\) be
    \(n\) times differentiable at an interior point \(x_0\) of \(I\). Suppose that
    \[
        P(x) = \sum_{k=0}^n a_k {(x-x_0)}^k
    \]
    is a polynomial of degree at most \(n\) such that
    \[
        f(x)=P(x)+\littleO\!\left({(x-x_0)}^n\right)
        \qquad\text{as }x\to x_0.
    \]
    Then \(P=P_n\), where \(P_n\) is the Taylor polynomial of \(f\) of order
    \(n\), centered at \(x_0\). In particular,
    \[
        a_k=\frac{f^{(k)}(x_0)}{k!}
        \qquad\text{for every }k\in\{0,\ldots,n\}.
    \]
\end{snippettheorem}

\begin{snippetproof}{taylor-uniqueness-proof}{taylor-uniqueness-theorem}{Uniqueness of the Taylor polynomial}
    Taylor's formula with the Peano remainder gives
    \[
        f(x)=P_n(x)+\littleO\!\left({(x-x_0)}^n\right).
    \]
    Hence
    \[
        P(x)-P_n(x)
        =[f(x)-P_n(x)]-[f(x)-P(x)]
        =\littleO\!\left({(x-x_0)}^n\right).
    \]
    Set \(Q=P-P_n\), and write
    \[
        Q(x) = \sum_{k=0}^n b_k{(x-x_0)}^k.
    \]
    Suppose, by contradiction, that \(Q\neq 0\), and let \(k_0\) be the
    smallest index such that \(b_{k_0}\neq 0\). Then
    \[
        \frac{Q(x)}{{(x-x_0)}^n}
        ={(x-x_0)}^{k_0-n}
        \sum_{k=k_0}^n b_k{(x-x_0)}^{k-k_0}.
    \]
    The sum tends to \(b_{k_0}\neq 0\). If \(k_0=n\), the quotient tends to
    \(b_n\neq 0\); if \(k_0<n\), its absolute value is unbounded. In either case
    the quotient cannot tend to zero, contradicting
    \(Q(x)=\littleO\!\left({(x-x_0)}^n\right)\). Therefore \(Q=0\), and hence
    \(P=P_n\).
\end{snippetproof}

\end{document}
